\documentclass{article}
\usepackage{amsthm,amsmath}
\newtheorem{definition}{Definition}
\newtheorem{theorem}{Theorem}

\begin{document}
\begin{definition}
Let $G$ and $H$ be groups with operations $+_G$ and $+_H$, respectively
For convenience, we shall use $+$ to denote either operation when
the intended operation is clear from context.

The direct product of $G$ and $H$,
denoted by $G \times H$, consists of all ordered pairs $(x,y)$
where $x \in G$ and $y \in H$. In other words,
\[
G \times H = \{(x,y) \mid x \in G \quad \text{and} \quad y \in H\}
\]

Also, let $+$ denote the operation of $G \times H$ as follows:
\[
(x,y) + (x',y') = (x+x',y+y')
\]
where $(x,y),(x',y') \in G \times H$.
\end{definition}

\begin{theorem}
If $G$ and $H$ are abelian groups, then $G \times H$ is also
an abelian group.
\end{theorem}

\begin{proof}
Let $G$ and $H$ be abelian groups, and let
$(x,y),(x',y'),(x'',y'') \in G \times H$.

\textbf{Commutativity:} 
\begin{align*}
(x,y)+(x',y')&=(x+x',y+y')\\
&=(x'+x,y'+y)\\
&=(x',y')+(x,y)
\end{align*}

Since $(x,y)+(x',y')=(x',y')+(x,y)$, it follows that
$G \times H$ is commutative.

\textbf{Associativity:}
\begin{align*}
(x,y)+((x',y')+(x'',y''))&=(x,y)+(x'+x'',y'+y'')\\
&=(x+(x'+x''),y+(y'+y''))\\
&=((x+x')+x'',(y+y')+y'')\\
&=((x,y)+(x',y'))+(x'',y'')
\end{align*}

Hence, $G \times H$ is associative.

\textbf{Identity:}

Let $(a,b) \in G \times H$. Consider the following equation
\[
(x,y)+(a,b)=(x,y)
\]

We have $(x+a,y+b)=(x,y)$. Hence,
\[
\begin{cases}
x+a=x\\
y+b=y
\end{cases}
\]

From the first equation, by adding $-x$ on both sides, we obtain
\begin{align*}
(x+a)+(-x)&=x+(-x)\\
a+(x+(-x))&=e_G\\
a+e_G&=e_G\\
a&=e_G
\end{align*}
where $e_G$ is the identity of $G$.

Similarly, from the second equation, we obtain $b=e_H$ where $e_H$
is the identity of $H$.

Now,
\begin{align*}
(x,y)+(e_G,e_H)&=(x+e_G,y+e_H)\\
&=(x,y)
\end{align*}

Since $G \times H$ is commutative, we also have
$(e_G,e_H)+(x,y)=(x,y)$

Hence, $G \times H$ has an identity element, namely $(e_G,e_H)$.

\textbf{Inverse:}
Let $(a,b) \in G \times H$. Consider the following equation
\[
(x,y)+(a,b)=(e_G,e_H)
\]

We have $x+a=e_G$. Hence, $a=-x$. Similarly, $b=-y$.

Now,
\begin{align*}
(x,y)+(-x,-y)&=(x+(-x),y+(-y))\\
&=(e_G,e_H)
\end{align*}

Since $G \times H$ is commutative, it is also the case that
$(-x,-y)+(x,y)=(e_G,e_H)$.

Furthermore, since $(e_G,e_H)$ is the identity of $G \times H$, it follows
that $(x,y)$ has an inverse, namely $(-x,-y)$.

Therefore, $G \times H$ is an abelian group.
\end{proof}
\end{document}
